
%% ====================== ปก ======================
\begin{frame}[plain]
  \maketitle
\end{frame}

%% ====================== เชื่อมจากบทที่ 1 ======================
\begin{frame}{จากบทที่ 1 (ตรรกศาสตร์) สู่บทที่ 2 (ทฤษฎีเซต)}
  บทที่ 1 เราเรียน ``ภาษาของการให้เหตุผล'' บทนี้เราเรียน ``ภาษาของการจัดกลุ่ม''
  \begin{center}
  \begin{tabular}{|l|l|}
    \hline
    \textbf{ตรรกศาสตร์ (บทที่ 1)} & \textbf{ทฤษฎีเซต (บทที่ 2)} \\ \hline
    ประพจน์ $p$ & เซต $A$ \\ \hline
    ``และ'' $\land$ & อินเตอร์เซกชัน $\cap$ \\ \hline
    ``หรือ'' $\lor$ & ยูเนียน $\cup$ \\ \hline
    นิเสธ $\neg$ & คอมพลีเมนต์ $^c$ \\ \hline
    สัจนิรันดร์ (T) & เอกภพสัมพัทธ์ $U$ \\ \hline
    ข้อขัดแย้ง (F) & เซตว่าง $\emptyset$ \\ \hline
  \end{tabular}
  \end{center}
  \vspace{0.3em}
  \begin{exampleblock}{}
    กฎทุกข้อที่เรียนในบทที่ 1 (เดอมอร์แกน แจกแจง ดูดกลืน ฯลฯ) จะกลับมาอีกครั้งใน ``เสื้อคลุมของเซต''
  \end{exampleblock}
\end{frame}

%% ====================== วัตถุประสงค์ ======================
\begin{frame}{วัตถุประสงค์การเรียนรู้ (สัปดาห์ที่ 4)}
  เมื่อเรียนจบสัปดาห์นี้ นักศึกษาจะสามารถ
  \begin{enumerate}
    \item อธิบายความหมายของ \textbf{เซตและสมาชิก} และใช้สัญลักษณ์ $\in,\ \notin$ ได้ถูกต้อง
    \item เขียนเซตแบบ \textbf{แจกแจงสมาชิก} และแบบ \textbf{บอกเงื่อนไข} และแปลงระหว่างสองแบบได้
    \item จำแนก \textbf{เซตว่าง เซตจำกัด เซตอนันต์} และเอกภพสัมพัทธ์ได้
    \item ตรวจสอบ \textbf{เซตย่อย} ($\subseteq$) \textbf{เซตย่อยแท้} ($\subsetneq$) และแยกความต่างระหว่าง $\in$ กับ $\subseteq$ ได้
  \end{enumerate}
  \vspace{0.5em}
  \begin{exampleblock}{สัปดาห์ที่ 5--6 จะต่อด้วย}
    การดำเนินการระหว่างเซต $\rightarrow$ กฎเดอมอร์แกน $\rightarrow$ เพาเวอร์เซต
    $\rightarrow$ ผลคูณคาร์ทีเซียน $\rightarrow$ การนับสมาชิกและการประยุกต์
  \end{exampleblock}
\end{frame}

%% ====================== หัวข้อ ======================
\begin{frame}{หัวข้อที่จะศึกษาในสัปดาห์นี้}
  \tableofcontents
\end{frame}


%% =====================================================================
\section{บทนำ: ทำไมต้องเรียนทฤษฎีเซต}
%% =====================================================================

\begin{frame}{เซตอยู่รอบตัวเรา}
  ในชีวิตประจำวัน เราจัดกลุ่มสิ่งต่าง ๆ อยู่ตลอดเวลา
  \begin{itemize}
    \item เพลย์ลิสต์เพลงโปรด = \textbf{กลุ่มของเพลง}
    \item รายชื่อเพื่อนในกลุ่ม LINE = \textbf{กลุ่มของคน}
    \item ตะกร้าสินค้าออนไลน์ = \textbf{กลุ่มของสินค้า}
    \item โฟลเดอร์ในเครื่อง = \textbf{กลุ่มของไฟล์}
  \end{itemize}
  \vspace{0.5em}
  \begin{block}{ในทางคณิตศาสตร์}
    กลุ่มของสิ่งต่าง ๆ เรียกว่า \textbf{เซต (Set)} และสิ่งที่อยู่ในกลุ่มเรียกว่า \textbf{สมาชิก (Element)}
  \end{block}
  \vspace{0.4em}
  ทฤษฎีเซตคือ ``ภาษากลาง'' ที่คณิตศาสตร์ทุกสาขาใช้ร่วมกัน เหมือนตัวอักษรที่เป็นพื้นฐานของภาษา
\end{frame}

\begin{frame}{ประวัติโดยสังเขป: Georg Cantor}
  \begin{itemize}
    \item ทฤษฎีเซตพัฒนาอย่างเป็นระบบโดย \textbf{Georg Cantor} (1845--1918)
          นักคณิตศาสตร์ชาวเยอรมัน ช่วงปลายศตวรรษที่ 19
    \item จุดเริ่มต้น: ศึกษาอนุกรมฟูเรียร์ แล้วตั้งคำถามกับธรรมชาติของ \textbf{ความไม่สิ้นสุด (Infinity)}
    \item ผลงานที่สั่นสะเทือนวงการ: \textbf{ความอนันต์มีหลายระดับ!}
      \begin{itemize}
        \item เซตจำนวนเต็ม $\mathbb{Z}$ กับจำนวนตรรกยะ $\mathbb{Q}$ มี ``ขนาดเท่ากัน''
        \item แต่เซตจำนวนจริง $\mathbb{R}$ ``ใหญ่กว่า'' อย่างพิสูจน์ได้
      \end{itemize}
  \end{itemize}
  \vspace{0.4em}
  \begin{exampleblock}{เกร็ด}
    แนวคิดของ Cantor ถูกต่อต้านอย่างหนักในยุคแรก แต่ปัจจุบันเป็น \textbf{รากฐานของคณิตศาสตร์ทั้งหมด}
    (เราจะได้ชิมเรื่อง ``อนันต์หลายระดับ'' ในสัปดาห์ที่ 6)
  \end{exampleblock}
\end{frame}

\begin{frame}{เซตกับวิทยาการคอมพิวเตอร์}
  \begin{itemize}
    \item \textbf{โครงสร้างข้อมูล} \eng{Data Structures} -- Array, List, Hash Set, Tree ล้วนคือการจัดเก็บ ``กลุ่มของข้อมูล''
    \item \textbf{ฐานข้อมูล} \eng{Database} -- SQL คือการดำเนินการระหว่างเซต: \texttt{UNION}, \texttt{INTERSECT}, \texttt{EXCEPT}
    \item \textbf{ทฤษฎีภาษา} \eng{Language Theory} -- ตัวอักษร สายอักขระ และภาษา ล้วนนิยามด้วยเซต
    \item \textbf{ความน่าจะเป็น} -- ปริภูมิตัวอย่างและเหตุการณ์ คือเซตและเซตย่อย
    \item \textbf{ปัญญาประดิษฐ์} -- ฐานความรู้และการอนุมาน คือการดำเนินการบนเซตของข้อเท็จจริง
  \end{itemize}
  \vspace{0.4em}
  \begin{exampleblock}{มุมมองสำคัญ}
    เข้าใจเซต = เข้าใจรากของ SQL, โครงสร้างข้อมูล และคณิตศาสตร์ทุกบทที่เหลือของวิชานี้
  \end{exampleblock}
\end{frame}

\begin{frame}{เซตในภาษา Python}
  ภาษาโปรแกรมสมัยใหม่มีชนิดข้อมูล ``เซต'' ให้ใช้โดยตรง
  \begin{block}{ตัวอย่าง Python}
    \texttt{>>> A = \{1, 2, 2, 3, 3, 3\}} \\
    \texttt{>>> A} \\
    \texttt{\{1, 2, 3\}} \hfill \eng{สมาชิกซ้ำถูกตัดทิ้งอัตโนมัติ} \\[0.4em]
    \texttt{>>> 2 in A} \hfill \eng{ทดสอบการเป็นสมาชิก} \\
    \texttt{True}
  \end{block}
  \vspace{0.4em}
  \begin{itemize}
    \item เซตของ Python ทำงานตรงตามนิยามคณิตศาสตร์: \textbf{ไม่มีลำดับ ไม่มีสมาชิกซ้ำ}
    \item ถ้าต้องการเก็บซ้ำ/มีลำดับ ต้องใช้ \texttt{list} แทน --- นี่คือเหตุผลที่ต้องแยกแนวคิดให้ชัด
  \end{itemize}
\end{frame}


%% =====================================================================
\section{นิยามพื้นฐาน: เซต สมาชิก และวิธีเขียนเซต}
%% =====================================================================

\begin{frame}{เซตและสมาชิก}
  \begin{block}{นิยาม}
    \textbf{เซต (Set)} คือ กลุ่มของวัตถุที่ระบุไว้อย่างชัดเจน เรียกวัตถุแต่ละชิ้นว่า
    \textbf{สมาชิก (Element)} ของเซต
    \begin{itemize}
      \item $a \in A$ อ่านว่า ``$a$ \textbf{เป็นสมาชิกของ} $A$''
      \item $a \notin A$ อ่านว่า ``$a$ \textbf{ไม่เป็นสมาชิกของ} $A$''
    \end{itemize}
  \end{block}
  \vspace{0.3em}
  ตัวอย่าง: ให้ $A = \{1, 2, 3, 4, 5\}$
  \begin{itemize}
    \item $3 \in A$ \hfill $\Rightarrow$ \TT
    \item $7 \in A$ \hfill $\Rightarrow$ \FF \quad (นั่นคือ $7 \notin A$)
  \end{itemize}
  \vspace{0.3em}
  ข้อความ $a \in A$ ตัดสินได้ว่าจริงหรือเท็จเสมอ จึงเป็น \textbf{ประพจน์} ตามบทที่ 1
  --- เขียนเป็นกฎกีดกันกลางได้ว่า $(a \in A) \lor (a \notin A)$
\end{frame}

\begin{frame}{เซตของจำนวนที่ต้องรู้จัก}
  \begin{columns}[T]
    \begin{column}{0.52\textwidth}
      \begin{itemize}\small
        \item $\mathbb{N} = \{0, 1, 2, 3, \ldots\}$ \\ \textbf{จำนวนนับ} \eng{Natural Numbers}
        \item $\mathbb{Z} = \{\ldots, -2, -1, 0, 1, 2, \ldots\}$ \\ \textbf{จำนวนเต็ม} \eng{Integers, มาจาก Zahlen}
        \item $\mathbb{Q} = \left\{\tfrac{a}{b} \,\middle|\, a, b \in \mathbb{Z},\ b \neq 0\right\}$ \\ \textbf{จำนวนตรรกยะ} \eng{Rational, มาจาก Quotient}
        \item $\mathbb{R}$ = \textbf{จำนวนจริง} \eng{Real Numbers} \\ รวมจำนวนอตรรกยะ เช่น $\sqrt{2}, \pi, e$
      \end{itemize}
    \end{column}
    \begin{column}{0.46\textwidth}
      \begin{center}
      \begin{tikzpicture}[scale=0.78,every node/.style={font=\small}]
        \draw[fill=lightnavy] (0,0) ellipse (3.1 and 1.95);
        \node at (0,1.62) {$\mathbb{R}$};
        \draw[fill=lyellow] (-0.25,-0.18) ellipse (2.3 and 1.35);
        \node at (-0.25,0.92) {$\mathbb{Q}$};
        \draw[fill=mint] (-0.5,-0.35) ellipse (1.55 and 0.85);
        \node at (-0.5,0.28) {$\mathbb{Z}$};
        \draw[fill=lpink] (-0.7,-0.5) ellipse (0.85 and 0.42);
        \node at (-0.7,-0.5) {$\mathbb{N}$};
        \node at (2.05,-1.3) {$\sqrt{2},\ \pi$};
      \end{tikzpicture}
      \end{center}
      \centering\small
      $\mathbb{N} \subseteq \mathbb{Z} \subseteq \mathbb{Q} \subseteq \mathbb{R}$
    \end{column}
  \end{columns}
  \vspace{0.3em}
  \footnotesize ข้อตกลงของตำราเล่มนี้: $\mathbb{N}$ เริ่มที่ 0
\end{frame}

\begin{frame}{เช็กความเข้าใจเร็ว: จริงหรือเท็จ?}
  พิจารณาข้อความต่อไปนี้ (คิดในใจ 1 นาที แล้วร่วมกันเฉลย)
  \begin{enumerate}
    \item $\sqrt{2} \in \mathbb{Q}$
    \item $0.75 \in \mathbb{Q}$
    \item $-3 \in \mathbb{N}$
    \item $0.\overline{3} \in \mathbb{Q}$
    \item $\pi \in \mathbb{R}$
  \end{enumerate}
  \pause
  \begin{exampleblock}{เฉลย}
    \footnotesize
    1) \FF\ ($\sqrt{2}$ เป็นจำนวนอตรรกยะ) \quad
    2) \TT\ ($0.75 = \tfrac{3}{4}$) \quad
    3) \FF\ ($\mathbb{N}$ ไม่มีจำนวนลบ แต่ $-3 \in \mathbb{Z}$) \quad
    4) \TT\ ($0.\overline{3} = \tfrac{1}{3}$) \quad
    5) \TT
  \end{exampleblock}
\end{frame}

\begin{frame}{วิธีเขียนเซตแบบที่ 1: แจกแจงสมาชิก}
  \begin{block}{แจกแจงสมาชิก (Roster Method)}
    เขียนสมาชิก \textbf{ทุกตัว} ในวงเล็บปีกกา $\{\ \}$ คั่นด้วยจุลภาค เช่น
    $A = \{1, 2, 3\}$ \quad $B = \{a, e, i, o, u\}$
  \end{block}
  \vspace{0.3em}
  กติกาสำคัญ 2 ข้อของเซต
  \begin{enumerate}
    \item \textbf{ไม่มีลำดับ} \eng{Unordered}: $\{1, 2\} = \{2, 1\}$
    \item \textbf{ไม่มีสมาชิกซ้ำ} \eng{No Duplicates}: $\{1, 1, 2, 2, 3\} = \{1, 2, 3\}$
  \end{enumerate}
  \vspace{0.3em}
  \begin{exampleblock}{เมื่อสมาชิกเยอะแต่มีแบบรูปชัดเจน ใช้จุดไข่ปลา (ellipsis)}
    $\{1, 2, 3, \ldots, 100\}$ = จำนวนเต็ม 1 ถึง 100 \\
    \footnotesize ข้อควรระวัง: เขียนสมาชิกอย่างน้อย 2--3 ตัวก่อน $\ldots$ เพื่อให้ผู้อ่านเดาแบบรูปได้
  \end{exampleblock}
\end{frame}

\begin{frame}{วิธีเขียนเซตแบบที่ 2: บอกเงื่อนไข}
  \begin{block}{บอกเงื่อนไขของสมาชิก (Set-builder Notation)}
    เขียนในรูป $\{x \mid \text{เงื่อนไขที่ } x \text{ ต้องเป็นจริง}\}$
    อ่านว่า ``เซตของ $x$ \textbf{โดยที่} \ldots'' เหมาะกับเซตที่มีสมาชิกมากหรือเป็นอนันต์
  \end{block}
  \vspace{0.3em}
  ตัวอย่าง
  \begin{itemize}
    \item $\{x \in \mathbb{Z} \mid x > 0\}$ = จำนวนเต็มบวก ($\mathbb{Z}^+$)
    \item $\{x \in \mathbb{Z} \mid x \text{ หารด้วย 2 ลงตัว}\}$ = จำนวนคู่
    \item $\{x \in \mathbb{R} \mid x^2 - 5x + 6 = 0\} = \{2, 3\}$ = เซตคำตอบของสมการ
  \end{itemize}
  \vspace{0.3em}
  \begin{exampleblock}{เทียบกับการเขียนโปรแกรม}
    Set-builder คล้าย list comprehension ของ Python: \\
    \texttt{\{x for x in range(100) if x \% 2 == 0\}} $\;\approx\; \{x \in \mathbb{N} \mid x < 100,\ x \text{ เป็นจำนวนคู่}\}$
  \end{exampleblock}
\end{frame}

\begin{frame}{ตัวอย่าง: แปลงจากบอกเงื่อนไขเป็นแจกแจงสมาชิก}
  จงเขียนเซตต่อไปนี้แบบแจกแจงสมาชิก
  \begin{enumerate}
    \item $B = \{n^2 \mid n \in \mathbb{N} \text{ และ } n < 5\}$
    \item $C = \{x \in \mathbb{R} \mid x^2 - 5x + 6 = 0\}$
    \item $D = \{x \in \mathbb{R} \mid x^2 + 1 = 0\}$
  \end{enumerate}
  \pause
  \begin{exampleblock}{เฉลย}
    \begin{enumerate}\footnotesize
      \item แทน $n = 0, 1, 2, 3, 4$ ได้ $B = \{0, 1, 4, 9, 16\}$
      \item แยกตัวประกอบ $(x-2)(x-3) = 0$ ได้ $C = \{2, 3\}$
      \item ไม่มีจำนวนจริงที่ $x^2 = -1$ ได้ $D = \emptyset$ \quad (เซตว่าง!)
    \end{enumerate}
  \end{exampleblock}
\end{frame}

\begin{frame}{เซตพิเศษที่ต้องรู้จัก}
  \begin{itemize}
    \item \textbf{เซตว่าง (Empty Set)} $\emptyset$ หรือ $\{\ \}$ --- เซตที่ไม่มีสมาชิกเลย
    \item \textbf{เซตที่มีสมาชิกตัวเดียว (Singleton Set)} เช่น $\{5\},\ \{a\},\ \{\emptyset\}$
    \item \textbf{เอกภพสัมพัทธ์ (Universal Set)} $U$ --- ``กรอบอ้างอิง'' ที่รวมทุกสิ่งที่สนใจในบริบทนั้น
          เช่น ศึกษาเรื่องจำนวน อาจให้ $U = \mathbb{R}$
    \item \textbf{เซตจำกัด (Finite Set)} --- นับจำนวนสมาชิกได้จบ เช่น $\{1,2,3\}$
    \item \textbf{เซตอนันต์ (Infinite Set)} --- สมาชิกไม่สิ้นสุด เช่น $\mathbb{N}, \mathbb{Z}, \mathbb{Q}, \mathbb{R}$
  \end{itemize}
  \vspace{0.4em}
  \begin{exampleblock}{}
    เซตว่างเป็นเซตย่อยของ \textbf{ทุกเซต}: $\emptyset \subseteq A$ เสมอ
    (จริงแบบ vacuous truth --- ``สมาชิกทุกตัวของ $\emptyset$'' ไม่มีให้ขัดแย้ง)
  \end{exampleblock}
\end{frame}

\begin{frame}{ระวัง! $\emptyset$ ไม่เท่ากับ $\{\emptyset\}$}
  \begin{alertblock}{ความเข้าใจผิดที่พบบ่อย}
    $\emptyset$ คือเซตที่ \textbf{ไม่มีสมาชิกเลย} \quad แต่ $\{\emptyset\}$ คือเซตที่ \textbf{มีสมาชิก 1 ตัว}
    (สมาชิกตัวนั้นคือ $\emptyset$)
  \end{alertblock}
  \vspace{0.4em}
  \begin{center}
  \begin{tabular}{|c|c|c|}
    \hline
    \textbf{เซต} & \textbf{จำนวนสมาชิก} & \textbf{เปรียบเทียบ} \\ \hline
    $\emptyset$ & 0 & กล่องเปล่า \\ \hline
    $\{\emptyset\}$ & 1 & กล่องที่มี ``กล่องเปล่า'' ข้างใน \\ \hline
    $\{\{\emptyset\}\}$ & 1 & กล่องที่มี ``กล่องซ้อนกล่องเปล่า'' \\ \hline
  \end{tabular}
  \end{center}
  \vspace{0.4em}
  ดังนั้น $\emptyset \neq \{\emptyset\} \neq \{\{\emptyset\}\}$ \quad และ $|\emptyset| = 0$ แต่ $|\{\emptyset\}| = 1$
\end{frame}


%% =====================================================================
\section{ความสัมพันธ์ระหว่างเซต: เซตย่อยและความเท่ากัน}
%% =====================================================================

\begin{frame}{เซตย่อย (Subset)}
  \begin{block}{นิยาม}
    $A$ เป็น \textbf{เซตย่อย} ของ $B$ เขียน $A \subseteq B$ ก็ต่อเมื่อ
    \textbf{ทุกสมาชิกของ $A$ เป็นสมาชิกของ $B$} ด้วย \\
    เขียนด้วยตัวบ่งปริมาณ (บทที่ 1): $\forall x\, (x \in A \rightarrow x \in B)$
  \end{block}
  \begin{columns}[T]
    \begin{column}{0.55\textwidth}
      ตัวอย่าง
      \begin{itemize}\small
        \item $\{1, 2\} \subseteq \{1, 2, 3\}$ \hfill \TT
        \item $\{1, 2\} \subseteq \{2, 3\}$ \hfill \FF\ \footnotesize(เพราะ $1 \notin \{2,3\}$)
        \item $\emptyset \subseteq \{1, 2\}$ \hfill \TT
        \item $A \subseteq A$ เสมอ \hfill \TT
      \end{itemize}
    \end{column}
    \begin{column}{0.4\textwidth}
      \begin{center}
      \begin{tikzpicture}[scale=0.72,every node/.style={font=\small}]
        \draw (-2.1,-1.5) rectangle (2.1,1.6);
        \draw[fill=lightnavy] (0,0) circle (1.25);
        \draw[fill=mint] (-0.3,-0.2) circle (0.62);
        \node at (0.75,0.85) {$B$};
        \node at (-0.3,-0.2) {$A$};
        \node at (1.85,1.38) {$U$};
      \end{tikzpicture}\\
      $A \subseteq B$
      \end{center}
    \end{column}
  \end{columns}
  \vspace{0.2em}
  \footnotesize วิธีพิสูจน์ $A \subseteq B$: สมมติ $x \in A$ แล้วแสดงว่า $x \in B$ (direct proof)
\end{frame}

\begin{frame}{เซตย่อยแท้ และความเท่ากันของเซต}
  \begin{block}{เซตย่อยแท้ (Proper Subset)}
    $A \subsetneq B$ ก็ต่อเมื่อ $A \subseteq B$ \textbf{และ} $A \neq B$
    (เป็นเซตย่อยแต่ ``เล็กกว่าจริง ๆ'') เช่น $\{1, 2\} \subsetneq \{1, 2, 3\}$
  \end{block}
  \vspace{0.4em}
  \begin{block}{ความเท่ากันของเซต (Extensionality)}
    \centering
    $A = B$ \quad ก็ต่อเมื่อ \quad $A \subseteq B$ \textbf{และ} $B \subseteq A$
  \end{block}
  \vspace{0.4em}
  \begin{exampleblock}{เทคนิคพิสูจน์มาตรฐาน: double inclusion}
    จะพิสูจน์ว่าเซตสองเซตเท่ากัน ให้พิสูจน์การเป็นเซตย่อย \textbf{สองทิศทาง} \\
    (1) $A \subseteq B$: ทุก $x \in A \Rightarrow x \in B$ \qquad
    (2) $B \subseteq A$: ทุก $x \in B \Rightarrow x \in A$ \\
    เทคนิคนี้จะใช้หนักมากในสัปดาห์ที่ 5!
  \end{exampleblock}
\end{frame}

\begin{frame}{จุดที่พลาดบ่อยที่สุด: $\in$ กับ $\subseteq$}
  \begin{alertblock}{ต่างกันอย่างไร}
    $\in$ = ความสัมพันธ์ระหว่าง \textbf{สมาชิก} กับ \textbf{เซต} \qquad
    $\subseteq$ = ความสัมพันธ์ระหว่าง \textbf{เซต} กับ \textbf{เซต}
  \end{alertblock}
  \vspace{0.3em}
  ให้ $A = \{1, 2, 3\}$
  \begin{center}
  \begin{tabular}{|l|c|l|}
    \hline
    \textbf{ข้อความ} & \textbf{จริง/เท็จ} & \textbf{เหตุผล} \\ \hline
    $1 \in A$ & \TT & 1 เป็นสมาชิกของ $A$ \\ \hline
    $\{1\} \in A$ & \FF & สมาชิกของ $A$ คือ $1,2,3$ ไม่ใช่เซต $\{1\}$ \\ \hline
    $\{1\} \subseteq A$ & \TT & สมาชิกทุกตัวของ $\{1\}$ อยู่ใน $A$ \\ \hline
    $\emptyset \subseteq A$ & \TT & เซตว่างเป็นเซตย่อยของทุกเซต \\ \hline
    $\emptyset \in A$ & \FF & $\emptyset$ ไม่ใช่สมาชิกตัวใดของ $A$ \\ \hline
  \end{tabular}
  \end{center}
  \vspace{0.3em}
  \footnotesize แต่ทั้งคู่จริงพร้อมกันได้ เช่น $\emptyset \in \{\emptyset\}$ และ $\emptyset \subseteq \{\emptyset\}$ ต่างก็จริง!
\end{frame}

\begin{frame}{เกร็ดปิดท้าย: ปฏิทรรศน์ของช่างตัดผม}
  \begin{block}{เรื่องเล่า (Bertrand Russell, 1901)}
    ช่างตัดผมประจำหมู่บ้านประกาศว่า ``ผมจะโกนหนวดให้ \textbf{เฉพาะ} คนที่ไม่โกนหนวดตัวเอง''
    \\ \centering \textbf{คำถาม: แล้วใครโกนหนวดให้ช่างตัดผม?}
  \end{block}
  \vspace{0.3em}
  \begin{itemize}
    \item ถ้าช่างโกนให้ตัวเอง $\rightarrow$ ผิดประกาศ (เขาโกนให้เฉพาะคนที่ \emph{ไม่} โกนเอง)
    \item ถ้าช่างไม่โกนให้ตัวเอง $\rightarrow$ ต้องโกนให้ตัวเอง (ตามประกาศ) --- ขัดแย้งทั้งสองทาง!
  \end{itemize}
  \vspace{0.3em}
  ในภาษาเซต: $R = \{x \mid x \notin x\}$ แล้วถามว่า $R \in R$ หรือไม่ ---
  คือ \textbf{ปฏิทรรศน์ของรัสเซลล์ (Russell's Paradox)}
  \vspace{0.3em}
  \begin{exampleblock}{บทเรียน}
    ``กลุ่มของอะไรก็ได้'' ไม่เป็นเซตเสมอไป --- จึงต้องมี \textbf{เอกภพสัมพัทธ์ $U$} คอยกำกับขอบเขต
  \end{exampleblock}
\end{frame}


%% =====================================================================
\section{กิจกรรมในชั้นเรียน (แบบฝึกหัดทำในห้อง)}
%% =====================================================================

\begin{frame}{วิธีทำกิจกรรม}
  \begin{block}{คำชี้แจง}
    ให้นักศึกษาทำกิจกรรมต่อไปนี้ \textbf{ลงในใบกิจกรรม (Worksheet)} ภายในเวลาที่กำหนด
    แล้วร่วมเฉลยและอภิปรายพร้อมกัน
  \end{block}
  \vspace{0.5em}
  \begin{itemize}
    \item ทำเดี่ยวหรือจับคู่ก็ได้ (ตามที่อาจารย์กำหนด)
    \item เมื่อหมดเวลา อาจารย์จะคลิกเพื่อ \textbf{เฉลย} ทีละข้อ
    \item อาสาสมัครออกมาอธิบายวิธีคิดหน้าชั้น
  \end{itemize}
  \vspace{0.5em}
  \begin{exampleblock}{เกณฑ์}
    เน้นกระบวนการคิดและการให้เหตุผล ไม่ใช่แค่คำตอบสุดท้าย
  \end{exampleblock}
\end{frame}

\begin{frame}{กิจกรรมที่ 1: เขียนเซตสองแบบ \eng{6 นาที}}
  จงเขียนเซตต่อไปนี้ \textbf{แบบแจกแจงสมาชิก}
  \begin{enumerate}
    \item เซตของจำนวนเต็มบวกที่น้อยกว่า 8
    \item $\{x \in \mathbb{Z} \mid x^2 = 9\}$
    \item เซตของสระในคำว่า \texttt{computer} (ตัวอักษรภาษาอังกฤษ)
  \end{enumerate}
  และจงเขียนเซตต่อไปนี้ \textbf{แบบบอกเงื่อนไข}
  \begin{enumerate}
    \setcounter{enumi}{3}
    \item $\{2, 4, 6, 8, 10, \ldots\}$
  \end{enumerate}
  \pause
  \begin{exampleblock}{เฉลย}
    \footnotesize
    1) $\{1, 2, 3, 4, 5, 6, 7\}$ \quad
    2) $\{-3, 3\}$ (อย่าลืมค่าลบ!) \quad
    3) $\{o, u, e\}$ (ไม่ซ้ำ ไม่มีลำดับ) \quad
    4) $\{x \in \mathbb{Z} \mid x > 0 \text{ และ } x \text{ หารด้วย 2 ลงตัว}\}$ หรือ $\{2n \mid n \in \mathbb{Z}^+\}$
  \end{exampleblock}
\end{frame}

\begin{frame}{กิจกรรมที่ 2: จริงหรือเท็จ \eng{6 นาที}}
  กำหนด $A = \{1, 2, \{3\}\}$ จงพิจารณาว่าข้อความแต่ละข้อจริงหรือเท็จ
  \begin{enumerate}
    \item $1 \in A$
    \item $3 \in A$
    \item $\{3\} \in A$
    \item $\{3\} \subseteq A$
    \item $\{1, 2\} \subseteq A$
    \item $\emptyset \subseteq A$
  \end{enumerate}
  \pause
  \begin{exampleblock}{เฉลย}
    \footnotesize
    1) \TT \quad
    2) \FF\ (สมาชิกคือ $1, 2, \{3\}$ --- ตัวเลข 3 เดี่ยว ๆ ไม่ใช่สมาชิก) \quad
    3) \TT\ (เซต $\{3\}$ เป็นสมาชิกตัวหนึ่ง) \quad
    4) \FF\ (จะเป็นเซตย่อยต้องมี $3 \in A$ ซึ่งไม่จริง) \quad
    5) \TT \quad
    6) \TT
  \end{exampleblock}
\end{frame}

\begin{frame}{กิจกรรมที่ 3: จำแนกประเภทเซต \eng{6 นาที}}
  \textbf{ก)} จงระบุว่าเซตใดเป็น เซตว่าง / เซตจำกัด / เซตอนันต์
  \begin{enumerate}
    \item $\{x \in \mathbb{N} \mid x < 100\}$
    \item $\{x \in \mathbb{R} \mid x^2 < 0\}$
    \item $\{x \in \mathbb{Z} \mid x \text{ หารด้วย 5 ลงตัว}\}$
  \end{enumerate}
  \textbf{ข)} จงเขียน \textbf{เซตย่อยทั้งหมด} ของ $\{a, b\}$
  \pause
  \begin{exampleblock}{เฉลย}
    \footnotesize
    ก) 1) เซตจำกัด (มี 100 ตัว: $0$ ถึง $99$) \quad
    2) เซตว่าง (กำลังสองของจำนวนจริงไม่ติดลบ) \quad
    3) เซตอนันต์ ($\{\ldots, -10, -5, 0, 5, 10, \ldots\}$) \\[0.3em]
    ข) $\emptyset,\ \{a\},\ \{b\},\ \{a, b\}$ --- มีทั้งหมด 4 ตัว
    \eng{สัปดาห์ที่ 6 จะรู้ว่าทำไมได้ $2^2 = 4$}
  \end{exampleblock}
\end{frame}


%% =====================================================================
\section{เกมท้ายคาบ: อยู่เซตไหน?}
%% =====================================================================

\begin{frame}{เกมท้ายคาบ: ``อยู่เซตไหน?'' \eng{ประมาณ 10 นาที}}
  \begin{block}{กติกา}
    \begin{enumerate}
      \item แบ่งทีมตามแถวที่นั่ง (แถวละ 1 ทีม)
      \item อาจารย์ฉายคำถาม \textbf{จริง/เท็จ} ทีละข้อ ให้เวลาข้อละ \textbf{15 วินาที}
      \item หมดเวลา ทุกทีมชูมือพร้อมกัน: \textcolor{Tgreen}{\textbf{มือขวา = จริง (T)}} \quad
            \textcolor{Fred}{\textbf{มือซ้าย = เท็จ (F)}}
      \item ทีมที่สมาชิกตอบถูก \textbf{เกินครึ่งทีม} ได้ 1 แต้ม --- ข้อสุดท้ายเป็น \textbf{ข้อทอง} ได้ 2 แต้ม
    \end{enumerate}
  \end{block}
  \vspace{0.4em}
  \begin{exampleblock}{}
    \centering ทีมที่แต้มสูงสุด = แชมป์เซตประจำสัปดาห์ \quad ห้ามเปิดสไลด์ย้อน ห้ามเปิดโพย!
  \end{exampleblock}
\end{frame}

\begin{frame}{เกม: รอบที่ 1--3}
  \begin{enumerate}
    \item[\textbf{1.}] {\large $\{1, 2, 3\} = \{3, 1, 2\}$}
    \pause
    \item[] \hfill \textcolor{Tgreen}{\textbf{จริง}} --- เซตไม่มีลำดับ
    \pause
    \item[\textbf{2.}] {\large $|\{1, 1, 2, 2, 2\}| = 5$}
    \pause
    \item[] \hfill \textcolor{Fred}{\textbf{เท็จ}} --- สมาชิกซ้ำนับครั้งเดียว $\{1, 2\}$ มี 2 ตัว
    \pause
    \item[\textbf{3.}] {\large $\emptyset \in \{1, 2, 3\}$}
    \pause
    \item[] \hfill \textcolor{Fred}{\textbf{เท็จ}} --- $\emptyset$ ไม่ใช่สมาชิกของเซตนี้ (แต่ $\emptyset \subseteq$ จริง!)
  \end{enumerate}
\end{frame}

\begin{frame}{เกม: รอบที่ 4--5 และข้อทอง}
  \begin{enumerate}
    \item[\textbf{4.}] {\large $0.5 \in \mathbb{Z}$}
    \pause
    \item[] \hfill \textcolor{Fred}{\textbf{เท็จ}} --- $0.5$ ไม่ใช่จำนวนเต็ม (แต่ $0.5 \in \mathbb{Q}$)
    \pause
    \item[\textbf{5.}] {\large $\{a\} \subseteq \{a, b\}$}
    \pause
    \item[] \hfill \textcolor{Tgreen}{\textbf{จริง}} --- สมาชิกทุกตัวของ $\{a\}$ อยู่ใน $\{a,b\}$
    \pause
    \item[\textbf{ทอง}] {\large $\{\emptyset\} = \emptyset$} \quad \eng{2 แต้ม!}
    \pause
    \item[] \hfill \textcolor{Fred}{\textbf{เท็จ}} --- ฝั่งซ้ายมีสมาชิก 1 ตัว ฝั่งขวาไม่มีเลย
  \end{enumerate}
  \pause
  \vspace{0.4em}
  \begin{exampleblock}{}
    \centering รวมแต้ม ประกาศแชมป์ แล้วพบกันสัปดาห์หน้า!
  \end{exampleblock}
\end{frame}


%% ====================== สรุป + แบบฝึกหัด ======================
\begin{frame}{สรุปสัปดาห์ที่ 4}
  \begin{itemize}
    \item \textbf{เซต} = กลุ่มของสมาชิกที่ระบุชัดเจน ใช้ $\in$ / $\notin$ บอกการเป็นสมาชิก
    \item \textbf{เขียนเซตได้ 2 แบบ}: แจกแจงสมาชิก $\{1,2,3\}$ และบอกเงื่อนไข $\{x \mid \ldots\}$
    \item กติกาเซต: \textbf{ไม่มีลำดับ ไม่มีสมาชิกซ้ำ}
    \item เซตพิเศษ: เซตว่าง $\emptyset$, เอกภพสัมพัทธ์ $U$, เซตจำกัด/อนันต์ และ $\emptyset \neq \{\emptyset\}$
    \item \textbf{เซตย่อย} $A \subseteq B$: ทุกสมาชิกของ $A$ อยู่ใน $B$ \quad
          $A = B \Leftrightarrow A \subseteq B$ และ $B \subseteq A$
    \item $\in$ (สมาชิก--เซต) ไม่เหมือน $\subseteq$ (เซต--เซต) --- จุดออกข้อสอบยอดฮิต
  \end{itemize}
  \begin{exampleblock}{}
    \centering สัปดาห์ถัดไป: การดำเนินการระหว่างเซต ($\cup\ \cap\ -\ ^c$) แผนภาพเวนน์ และกฎเดอมอร์แกน
  \end{exampleblock}
\end{frame}

\begin{frame}{แบบฝึกหัดมอบหมาย (สัปดาห์ที่ 4)}
  ทำลงสมุด ส่งต้นคาบสัปดาห์หน้า (ข้อละ 2 คะแนน รวม 10 คะแนน)
  \begin{enumerate}\small
    \item เขียนเซตต่อไปนี้ทั้งแบบแจกแจงสมาชิกและแบบบอกเงื่อนไข:
          (ก) จำนวนเฉพาะที่น้อยกว่า 20 \quad (ข) จำนวนเต็ม $x$ ที่ $x^2 - 4 = 0$
    \item พิจารณาจริง/เท็จ พร้อมเหตุผล:
          (ก) $\emptyset \in \{\emptyset\}$ \quad (ข) $\emptyset \subseteq \{\emptyset\}$ \quad
          (ค) $\{1, 2\} \in \{\{1, 2\}\}$ \quad (ง) $\{1, 2\} \subseteq \{\{1\}, \{2\}\}$
    \item อธิบายความแตกต่างระหว่าง $\in$ และ $\subseteq$ พร้อมยกตัวอย่างของตนเอง 2 ตัวอย่าง
    \item เขียนเซตย่อย \textbf{ทั้งหมด} ของ $\{1, 2, 3\}$ (บอกจำนวนรวมด้วย)
    \item \textbf{(ประยุกต์)} ให้ $U$ = นักศึกษาทั้งหมดในห้องนี้, $A$ = คนที่มีโทรศัพท์ Android,
          $B$ = คนที่เคยเขียน Python จงบรรยายด้วยคำพูดว่าเซต
          $\{x \in U \mid x \in A \text{ และ } x \in B\}$ คือใคร และยกตัวอย่างสมาชิกถ้ามี
  \end{enumerate}
\end{frame}

\begin{frame}[plain]
  \vfill
  \centering
  {\Large\color{navy}\textbf{จบสัปดาห์ที่ 4}}\par
  \vspace{0.8em}
  {\color{gray} สัปดาห์ถัดไป: สัปดาห์ที่ 5 · การดำเนินการระหว่างเซต และกฎเดอมอร์แกน}
  \vfill
\end{frame}
